\par\Needspace{6\baselineskip}\medskip
\SourceLeaf{196}\begin{center}\hypertarget{native-196-article145}{}\pdfbookmark[2]{145.}{article-145}145.\end{center}
Postquam hoc modo corporis coelestis positiones in orbita determinatae sunt,
duplex elementorum calculus tum e combinatione loci secundi cum tertio,
tum e combinatione primi cum secundo, una cum temporum intervallis respondentibus,
inchoabitur. Antequam vero haec operatio suscipiatur, ipsa temporum
intervalla quadam correctione opus habent, siquidem constitutum fuerit,
secundum methodum tertiam art. 118 aberrationis rationem habere. In hocce
scilicet casu pro temporibus veris ficta substituenda sunt, illis resp. $493\rho,493\rho',493\rho''$
minutis secundis anteriora. Pro computandis distantiis $\rho,\rho',\rho''$ habemus
formulas
\[
\begin{aligned}
\rho&=\frac{R\sin(AD'-\zeta)}{\sin(\zeta-AD'+\delta)}=\frac{r\sin(AD'-\zeta)}{\sin\delta}\\[3pt]
\rho'&=\frac{R'\sin(\delta'-z)}{\sin z}=\frac{r'\sin(\delta'-z)}{\sin\delta'}\\[3pt]
\rho''&=\frac{R''\sin(A''D'-\zeta'')}{\sin(\zeta''-A''D'+\delta'')}=\frac{r''\sin(A''D'-\zeta'')}{\sin\delta''}\mathpunct{\cdot}
\end{aligned}
\]

Ceterum si observationes ab initio statim per methodum primam vel secundam
art. 118 ab aberratione purgatae fuissent, hicce calculus omittendus,
neque adeo necessarium foret, valores distantiarum $\rho,\rho',\rho''$ eruere, nisi forte
ad confirmandum, an ii, quibus calculus aberrationum superstructus erat, satis
exacti fuerint. Denique sponte patet, totum istum calculum tunc quoque supprimendum
esse, quando aberrationem omnino negligere placuerit.
